%%%%%%%%%%%%%%%%%%%%%%%%%%%%%%%%%%%%%%%%%%%%%%%%%%%%%%%%%%%%%%%%%%%%%%%%%%%%%%%%
% Document LaTeX : Équations différentielles
% Auteur : Laurent
% Date : 2026-07-11
% Description : Exemples d'équations différentielles et leurs solutions
%%%%%%%%%%%%%%%%%%%%%%%%%%%%%%%%%%%%%%%%%%%%%%%%%%%%%%%%%%%%%%%%%%%%%%%%%%%%%%%

\documentclass[12pt]{article}
\usepackage[utf8]{inputenc}
\usepackage[french]{babel}
\usepackage{amsmath}
\usepackage{amssymb}
\usepackage{amsthm}
\usepackage{graphicx}
\usepackage{hyperref}
\usepackage{geometry}
\usepackage{enumitem}
\usepackage{tcolorbox}
\usepackage{pgfplots}
\pgfplotsset{compat=1.18}

% Configuration de la page
\geometry{a4paper, margin=1in}

% Définition des environnements personnalisés
\newtheorem{exercice}{Exercice}
\newtheorem{solution}{Solution}
\newtheorem{theorem}{Théorème}[section]
\newtheorem{example}{Exemple}[section]

% Boîtes colorées pour les équations et solutions
\newtcolorbox{myequation}[1][]{colback=blue!5!white, colframe=blue!75!black, title=#1}
\newtcolorbox{mybox}[1][]{colback=green!5!white, colframe=green!75!black, title=#1}

% Titre et métadonnées
\title{Exemples d'Équations Différentielles}
\author{Laurent}
\date{\today}

%%%%%%%%%%%%%%%%%%%%%%%%%%%%%%%%%%%%%%%%%%%%%%%%%%%%%%%%%%%%%%%%%%%%%%%%%%%%%%%
\begin{document}

\maketitle

\tableofcontents

\newpage

%%%%%%%%%%%%%%%%%%%%%%%%%%%%%%%%%%%%%%%%%%%%%%%%%%%%%%%%%%%%%%%%%%%%%%%%%%%%%%%
\section{Introduction}
Les équations différentielles sont des équations qui relient une fonction à ses dérivées. Elles jouent un rôle fondamental en mathématiques, en physique, en ingénierie et dans de nombreuses autres disciplines scientifiques. Ce document présente divers types d'équations différentielles avec leurs solutions.

%%%%%%%%%%%%%%%%%%%%%%%%%%%%%%%%%%%%%%%%%%%%%%%%%%%%%%%%%%%%%%%%%%%%%%%%%%%%%%%
\section{Équations Différentielles du 1er Ordre}

\subsection{Équations à Variables Séparables}
\begin{myequation}[Forme générale]
Une équation différentielle à variables séparables a la forme :
\[
\frac{dy}{dx} = f(x) \cdot g(y)
\]
\end{myequation}

\begin{example}
Résoudre l'équation différentielle :
\[
\frac{dy}{dx} = 2xy
\]
\end{example}

\begin{mybox}[Solution]
On sépare les variables :
\[
\frac{dy}{y} = 2x \, dx
\]
On intègre les deux côtés :
\[
\int \frac{1}{y} \, dy = \int 2x \, dx \implies \ln|y| = x^2 + C
\]
On exponentie pour obtenir la solution générale :
\[
y = \pm e^{x^2 + C} = \pm e^C e^{x^2}
\]
En posant $C' = \pm e^C$, on obtient :
\[
\boxed{y = C' e^{x^2}}
\]
\end{mybox}

\begin{example}
Résoudre l'équation différentielle avec la condition initiale $y(0) = 3$ :
\[
\frac{dy}{dx} = 2xy
\]
\end{example}

\begin{mybox}[Solution]
D'après l'exemple précédent, la solution générale est $y = C e^{x^2}$. Avec la condition initiale $y(0) = 3$ :
\[
3 = C e^{0} \implies C = 3
\]
La solution est donc :
\[
\boxed{y = 3 e^{x^2}}
\]
\end{mybox}

%%%%%%%%%%%%%%%%%%%%%%%%%%%%%%%%%%%%%%%%%%%%%%%%%%%%%%%%%%%%%%%%%%%%%%%%%%%%%%%
\subsection{Équations Linéaires du 1er Ordre}
\begin{myequation}[Forme générale]
Une équation différentielle linéaire du 1er ordre a la forme :
\[
y' + a(x) y = b(x)
\]
\end{myequation}

\begin{theorem}[Solution générale]
La solution générale de l'équation $y' + a(x) y = b(x)$ est donnée par :
\[
y = \frac{1}{\mu(x)} \left( \int \mu(x) b(x) \, dx + C \right), \quad \text{où} \quad \mu(x) = e^{\int a(x) \, dx}
\]
\end{theorem}

\begin{example}
Résoudre l'équation différentielle :
\[
y' + 3y = e^{-2x}
\]
\end{example}

\begin{mybox}[Solution]
Le facteur intégrant est :
\[
\mu(x) = e^{\int 3 \, dx} = e^{3x}
\]
On multiplie l'équation par $\mu(x)$ :
\[
e^{3x} y' + 3 e^{3x} y = e^{x}
\]
Le côté gauche est la dérivée de $e^{3x} y$ :
\[
\frac{d}{dx} \left( e^{3x} y \right) = e^{x}
\]
On intègre :
\[
e^{3x} y = e^{x} + C \implies y = e^{-2x} + C e^{-3x}
\]
La solution générale est donc :
\[
\boxed{y = e^{-2x} + C e^{-3x}}
\]
\end{mybox}

\begin{example}
Résoudre l'équation différentielle avec la condition initiale $y(0) = 1$ :
\[
y' + 3y = e^{-2x}
\]
\end{example}

\begin{mybox}[Solution]
D'après l'exemple précédent, la solution générale est $y = e^{-2x} + C e^{-3x}$. Avec la condition initiale $y(0) = 1$ :
\[
1 = e^{0} + C e^{0} \implies 1 = 1 + C \implies C = 0
\]
La solution est donc :
\[
\boxed{y = e^{-2x}}
\]
\end{mybox}

%%%%%%%%%%%%%%%%%%%%%%%%%%%%%%%%%%%%%%%%%%%%%%%%%%%%%%%%%%%%%%%%%%%%%%%%%%%%%%%
\subsection{Équations de Bernoulli}
\begin{myequation}[Forme générale]
Une équation différentielle de Bernoulli a la forme :
\[
y' + a(x) y = b(x) y^n
\]
\end{myequation}

\begin{example}
Résoudre l'équation différentielle :
\[
y' + y = y^3
\]
\end{example}

\begin{mybox}[Solution]
On pose $z = y^{1-n} = y^{-2}$, donc $z' = -2 y^{-3} y'$. L'équation devient :
\[
\frac{dz}{dx} = -2 y^{-3} y' = -2 (y^3 - y) y^{-3} = -2 (1 - y^{-2}) = -2 + 2z
\]
On a donc une équation linéaire en $z$ :
\[
z' - 2z = -2
\]
Le facteur intégrant est $\mu(x) = e^{\int -2 \, dx} = e^{-2x}$. On multiplie par $\mu(x)$ :
\[
e^{-2x} z' - 2 e^{-2x} z = -2 e^{-2x}
\]
Le côté gauche est la dérivée de $e^{-2x} z$ :
\[
\frac{d}{dx} \left( e^{-2x} z \right) = -2 e^{-2x}
\]
On intègre :
\[
e^{-2x} z = e^{-2x} + C \implies z = 1 + C e^{2x}
\]
On revient à $y$ :
\[
y^{-2} = 1 + C e^{2x} \implies y^2 = \frac{1}{1 + C e^{2x}} \implies y = \pm \frac{1}{\sqrt{1 + C e^{2x}}}
\]
La solution générale est donc :
\[
\boxed{y = \frac{1}{\sqrt{1 + C e^{2x}}}}
\]
\end{mybox}

%%%%%%%%%%%%%%%%%%%%%%%%%%%%%%%%%%%%%%%%%%%%%%%%%%%%%%%%%%%%%%%%%%%%%%%%%%%%%%%
\section{Équations Différentielles du 2nd Ordre}

\subsection{Équations Linéaires à Coefficients Constants}
\begin{myequation}[Forme générale]
Une équation différentielle linéaire du 2nd ordre à coefficients constants a la forme :
\[
a y'' + b y' + c y = f(x)
\]
\end{myequation}

\begin{theorem}[Solution de l'équation homogène]
Pour l'équation homogène $a y'' + b y' + c y = 0$, on cherche les racines de l'équation caractéristique :
\[
a r^2 + b r + c = 0
\]
Les solutions sont :
\begin{itemize}
    \item Si les racines $r_1$ et $r_2$ sont réelles et distinctes :
    \[
    y_h = C_1 e^{r_1 x} + C_2 e^{r_2 x}
    \]
    \item Si les racines sont réelles et égales ($r_1 = r_2$) :
    \[
    y_h = (C_1 + C_2 x) e^{r_1 x}
    \]
    \item Si les racines sont complexes ($r = \alpha \pm i \beta$) :
    \[
    y_h = e^{\alpha x} \left( C_1 \cos(\beta x) + C_2 \sin(\beta x) \right)
    \]
\end{itemize}
\end{theorem}

\begin{example}
Résoudre l'équation différentielle :
\[
y'' - 4y' + 4y = 0
\]
\end{example}

\begin{mybox}[Solution]
L'équation caractéristique est :
\[
r^2 - 4r + 4 = 0 \implies (r - 2)^2 = 0 \implies r = 2 \text{ (racine double)}
\]
La solution générale est donc :
\[
\boxed{y = (C_1 + C_2 x) e^{2x}}
\]
\end{mybox}

\begin{example}
Résoudre l'équation différentielle :
\[
y'' + y = 0
\]
\end{example}

\begin{mybox}[Solution]
L'équation caractéristique est :
\[
r^2 + 1 = 0 \implies r = \pm i
\]
La solution générale est donc :
\[
\boxed{y = C_1 \cos(x) + C_2 \sin(x)}
\]
\end{mybox}

\begin{example}
Résoudre l'équation différentielle :
\[
y'' + 2y' + 5y = 0
\]
\end{example}

\begin{mybox}[Solution]
L'équation caractéristique est :
\[
r^2 + 2r + 5 = 0 \implies r = \frac{-2 \pm \sqrt{4 - 20}}{2} = -1 \pm 2i
\]
La solution générale est donc :
\[
\boxed{y = e^{-x} \left( C_1 \cos(2x) + C_2 \sin(2x) \right)}
\]
\end{mybox}

%%%%%%%%%%%%%%%%%%%%%%%%%%%%%%%%%%%%%%%%%%%%%%%%%%%%%%%%%%%%%%%%%%%%%%%%%%%%%%%
\subsection{Équations Non Homogènes}
\begin{theorem}[Méthode des coefficients indéterminés]
Pour résoudre une équation non homogène $a y'' + b y' + c y = f(x)$, on cherche une solution particulière $y_p$ en fonction de la forme de $f(x)$ :
\begin{itemize}
    \item Si $f(x) = P_n(x)$ (polynôme de degré $n$), on essaie $y_p = Q_n(x)$ (polynôme de degré $n$).
    \item Si $f(x) = P_n(x) e^{kx}$, on essaie $y_p = Q_n(x) e^{kx}$.
    \item Si $f(x) = P_n(x) \cos(kx)$ ou $P_n(x) \sin(kx)$, on essaie $y_p = Q_n(x) \cos(kx) + R_n(x) \sin(kx)$.
\end{itemize}
Si $y_p$ est solution de l'équation homogène, on multiplie par $x$ ou $x^2$ selon la multiplicité.
\end{theorem}

\begin{example}
Résoudre l'équation différentielle :
\[
y'' + y = \sin(2x)
\]
\end{example}

\begin{mybox}[Solution]
L'équation homogène associée est $y'' + y = 0$, dont la solution générale est $y_h = C_1 \cos(x) + C_2 \sin(x)$. On cherche une solution particulière de la forme :
\[
y_p = A \cos(2x) + B \sin(2x)
\]
On calcule les dérivées :
\[
y_p' = -2A \sin(2x) + 2B \cos(2x)
\]
\[
y_p'' = -4A \cos(2x) - 4B \sin(2x)
\]
On substitue dans l'équation :
\[
y_p'' + y_p = (-4A + A) \cos(2x) + (-4B + B) \sin(2x) = -3A \cos(2x) - 3B \sin(2x) = \sin(2x)
\]
On identifie les coefficients :
\[
-3A = 0 \implies A = 0
\]
\[
-3B = 1 \implies B = -\frac{1}{3}
\]
La solution particulière est donc :
\[
y_p = -\frac{1}{3} \sin(2x)
\]
La solution générale est :
\[
\boxed{y = C_1 \cos(x) + C_2 \sin(x) - \frac{1}{3} \sin(2x)}
\]
\end{mybox}

%%%%%%%%%%%%%%%%%%%%%%%%%%%%%%%%%%%%%%%%%%%%%%%%%%%%%%%%%%%%%%%%%%%%%%%%%%%%%%%
\section{Systèmes d'Équations Différentielles}

\subsection{Systèmes Linéaires 2x2}
\begin{myequation}[Forme générale]
Un système linéaire 2x2 a la forme :
\[
\begin{cases}
    x' = a x + b y \\
    y' = c x + d y
\end{cases}
\]
\end{myequation}

\begin{theorem}[Solution générale]
La solution générale du système $X' = A X$ est donnée par :
\[
X = C_1 e^{\lambda_1 t} \mathbf{v}_1 + C_2 e^{\lambda_2 t} \mathbf{v}_2
\]
oindent où $\lambda_1$ et $\lambda_2$ sont les valeurs propres de $A$ et $\mathbf{v}_1$, $\mathbf{v}_2$ les vecteurs propres associés.
\end{theorem}

\begin{example}
Résoudre le système :
\[
\begin{cases}
    x' = 2x + y \\
    y' = x + 2y
\end{cases}
\]
\end{example}

\begin{mybox}[Solution]
On écrit le système sous forme matricielle :
\[
\begin{pmatrix} x \\ y \end{pmatrix}' = \begin{pmatrix} 2 & 1 \\ 1 & 2 \end{pmatrix} \begin{pmatrix} x \\ y \end{pmatrix}
\]
On cherche les valeurs propres de la matrice $A = \begin{pmatrix} 2 & 1 \\ 1 & 2 \end{pmatrix}$ :
\[
\det(A - \lambda I) = \begin{vmatrix} 2 - \lambda & 1 \\ 1 & 2 - \lambda \end{vmatrix} = (2 - \lambda)^2 - 1 = \lambda^2 - 4\lambda + 3 = 0
\]
Les racines sont :
\[
\lambda = \frac{4 \pm \sqrt{16 - 12}}{2} = \frac{4 \pm 2}{2} \implies \lambda_1 = 3, \lambda_2 = 1
\]
On cherche les vecteurs propres associés :
\begin{itemize}
    \item Pour $\lambda_1 = 3$ :
    \[
    (A - 3I) \mathbf{v}_1 = 0 \implies \begin{pmatrix} -1 & 1 \\ 1 & -1 \end{pmatrix} \begin{pmatrix} v_{11} \\ v_{12} \end{pmatrix} = 0 \implies v_{11} = v_{12}
    \]
    On choisit $\mathbf{v}_1 = \begin{pmatrix} 1 \\ 1 \end{pmatrix}$.
    \item Pour $\lambda_2 = 1$ :
    \[
    (A - I) \mathbf{v}_2 = 0 \implies \begin{pmatrix} 1 & 1 \\ 1 & 1 \end{pmatrix} \begin{pmatrix} v_{21} \\ v_{22} \end{pmatrix} = 0 \implies v_{21} = -v_{22}
    \]
    On choisit $\mathbf{v}_2 = \begin{pmatrix} 1 \\ -1 \end{pmatrix}$.
\end{itemize}
La solution générale est donc :
\[
\boxed{\begin{pmatrix} x \\ y \end{pmatrix} = C_1 e^{3t} \begin{pmatrix} 1 \\ 1 \end{pmatrix} + C_2 e^{t} \begin{pmatrix} 1 \\ -1 \end{pmatrix}}
\]
\end{mybox}

%%%%%%%%%%%%%%%%%%%%%%%%%%%%%%%%%%%%%%%%%%%%%%%%%%%%%%%%%%%%%%%%%%%%%%%%%%%%%%%
\section{Graphiques des Solutions}

\subsection{Solution de $y' = -y$}
L'équation $y' = -y$ a pour solution générale $y = C e^{-x}$. Voici le graphique pour différentes valeurs de $C$ :

\begin{tikzpicture}
\begin{axis}[
    axis lines = middle,
    xlabel = {$x$},
    ylabel = {$y$},
    xmin = -2, xmax = 2,
    ymin = -2, ymax = 2,
    grid = both,
    legend pos = north east
]

% Solutions pour C = 1, 2, -1
\addplot[blue, domain=-2:2, samples=100] {exp(-x)};
\addlegendentry{$C = 1$}

\addplot[red, domain=-2:2, samples=100] {2*exp(-x)};
\addlegendentry{$C = 2$}

\addplot[green, domain=-2:2, samples=100] {-exp(-x)};
\addlegendentry{$C = -1$}

\end{axis}
\end{tikzpicture}

\subsection{Solution de $y'' + y = 0$}
L'équation $y'' + y = 0$ a pour solution générale $y = A \cos(x) + B \sin(x)$. Voici le graphique pour $A = 1$ et $B = 0$ :

\begin{tikzpicture}
\begin{axis}[
    axis lines = middle,
    xlabel = {$x$},
    ylabel = {$y$},
    xmin = -2*pi, xmax = 2*pi,
    ymin = -2, ymax = 2,
    grid = both,
    trig format = rad
]

\addplot[blue, domain=-2*pi:2*pi, samples=200] {cos(deg(x))};
\addlegendentry{$y = \cos(x)$}

\end{axis}
\end{tikzpicture}

%%%%%%%%%%%%%%%%%%%%%%%%%%%%%%%%%%%%%%%%%%%%%%%%%%%%%%%%%%%%%%%%%%%%%%%%%%%%%%%
\section{Exercices Proposés}

\begin{exercice}
Résoudre l'équation différentielle suivante :
\[
\frac{dy}{dx} = x y^2
\]
\end{exercice}

\begin{exercice}
Résoudre l'équation différentielle suivante avec la condition initiale $y(0) = 2$ :
\[
y' + 2y = 3
\]
\end{exercice}

\begin{exercice}
Résoudre l'équation différentielle suivante :
\[
y'' - 5y' + 6y = 0
\]
\end{exercice}

\begin{exercice}
Résoudre l'équation différentielle suivante :
\[
y'' + 4y = \cos(2x)
\]
\end{exercice}

\begin{exercice}
Résoudre le système suivant :
\[
\begin{cases}
    x' = x + 2y \\
    y' = 3x + 2y
\end{cases}
\]
\end{exercice}

%%%%%%%%%%%%%%%%%%%%%%%%%%%%%%%%%%%%%%%%%%%%%%%%%%%%%%%%%%%%%%%%%%%%%%%%%%%%%%%
\section{Conclusion}
Les équations différentielles sont un outil puissant pour modéliser de nombreux phénomènes naturels et techniques. Ce document a présenté quelques types d'équations différentielles et leurs méthodes de résolution. Pour aller plus loin, tu peux explorer les équations aux dérivées partielles, les équations différentielles stochastiques, ou les méthodes numériques de résolution.

%%%%%%%%%%%%%%%%%%%%%%%%%%%%%%%%%%%%%%%%%%%%%%%%%%%%%%%%%%%%%%%%%%%%%%%%%%%%%%%
\end{document}
